As eqs. (5.35) e (5.38) estão relacionadas porque (5.35) é a equação diferencial de crescimento logístico e (5.38) é sua solução explícita obtida por separação de variáveis.

\textbf{A equação diferencial logística (5.35):}
\[
\frac{dN}{dt} = \lambda N\!\left(1 - \frac{N}{K}\right),
\]
onde $N$ é a população, $\lambda > 0$ a taxa de crescimento intrínseca e $K$ a capacidade de suporte.

\textbf{Interpretação:}
\begin{itemize}
\item Quando $N \ll K$: o fator $(1-N/K) \approx 1$ e $N$ cresce exponencialmente ($dN/dt \approx \lambda N$);
\item Quando $N \to K$: o fator $(1-N/K) \to 0$ e o crescimento cessa;
\item O equilíbrio estável ocorre em $N^* = K$.
\end{itemize}

\textbf{Relação com a eq. (5.38):}
A eq. (5.38) representa a população em função do tempo, integrando (5.35) por separação de variáveis:
\[
\int\frac{dN}{N(1-N/K)} = \int \lambda\,dt.
\]
O lado esquerdo é resolvido por frações parciais e conduz à solução logística explícita da eq. (5.38). Portanto, as duas equações descrevem a \textbf{mesma dinâmica} em formas complementares: diferencial e integral. 